% chktex-file 18
\section{Algebraic varieties}
\subsection{Reminder on algebraic geometry}
As the theory of algebraic groups is solely built upon that of algebraic geometry, it is worth recalling fundamental notions of that field. 

\subsubsection{The affine space}

\begin{defn}
    The \emph{affine $n$-space} is the set $\ank := K^n$.

    An \emph{affine variety} in $\ank$ is the vanishing set of an ideal in $K\left[x_1, \ldots, x_n\right]$ for some $n \geq 1$, i.e.\ it is a set $X \in \ank$ such that
    \[ X = V(I) := \left\{x \in \ank \ \middle| \  f(x) = 0 \text{ for all } f\in I\right\} \]
    for some ideal $I \in K[x_1, \ldots k_n]$.

    To an affine variety $X \subseteq \ank$, one may assign an ideal in $K\left[x_1, \ldots, x_n\right]$ of all polynomials vanishing on $X$
    \[I(X) := \left\{f \in K\left[x_1, \ldots, x_n\right] \ \middle| \  f(x) = 0 \text{ for all } x\in X\right\} \]
\end{defn}


Note that, since $K\left[x_1, \ldots, x_n\right]$ is noetherian, an affine variety may also be defined as the vanishing set of a finite number of polynomials. Recall that such sets define the closed sets of a topology on $\ank$, called the \emph{Zariski topology}. Thus classical notions of topology apply to affine varieties (or, more generally, to algebraic varieties). 

Notions of irreducibility and connectedness will be of particular interest later on. A topological space is called \emph{irreducible} if it cannot be written as the union of two distinct, non-empty, closed strict subsets. Similarly, a topological space is called \emph{connected} if it cannot be written as the union of two disjoint non-empty open subsets. \\
Irreducibility is a property that is characterized very explicitly for algebraic varieties, and irreducible varieties are exactly the building blocks of varieties.

\begin{theorem}
    Let $X \subseteq \ank$ be an affine variety. Then $X$ is irreducible if and only if $I(X) \subseteq K[x_1, \ldots, x_n]$ is prime.

    Moreover, $X$ can be decomposed uniquely into a union of irreducible components ${X = X_1 \cup \ldots \cup X_n}$ where no $X_i$ is contained in another.
\end{theorem}

\subsubsection{The projective space}

A similar story can be told about varieties in projective space. We recall the fundamental notions in that case.

\begin{defn}
The \emph{projective $n$-space} $\pnk$ is the set of all lines in $\affn{n+1}{K}$, or equivalently, the set of equivalence classes in $K^{n+1} \setminus \left\{(0,\ldots,0)\right\}$ relative to the equivalence
\[ (x_0, \ldots, x_n) \sim (y_0, \ldots, y_n)\]
if and only if there is $\lambda \in K^*$ such that $y_i = \lambda \cdot x_i$ for all $i$. We denote $[x_0:\ldots: x_n]$ the equivalence class of a point in $K^{n+1}\setminus\left\{\left(0, \ldots, 0\right)\right\}$ and call them \emph{homogeneous coordinates}.
\end{defn}

As projective space ignores scalar multiplication, not every polynomial can be seen as a function on projective space. This leads to the notion of homogeneity, for both polynomials and ideals.

\begin{defn}
    Let $f \in K[x_0, \ldots, x_n]$ be a polynomial. We can write $f = f_0 + f_1 + \cdots + f_d$, where $f_i$ contains all terms of degree $i$ and $d = \deg{f}$. $f$ is called \emph{homogeneous} if $f = f_d$, i.e.\ if all its monomials are of the same degree.

    An ideal $I \in K[x_0, \ldots, x_n]$ is \emph{homogenous} if it is generated by a set of homogeneous polynomials.
\end{defn}

Let $f \in K[x_0, \ldots, x_n]$ be a homogeneous polynomial of degree $d$ and $ {a = (a_0, \ldots, a_n) \in K^{n+1}}$. Since $f$ is homogeneous, for any $\lambda \in K$, we see that ${f(\lambda \cdot a) = \lambda^d f(a)}$ due to $f$ being homogeneous, so that $f$ can be seen as a map on projective space. This motivates the following definition.

\begin{defn}
    A \emph{projective variety} in $\pnk$ is the vanishing set of a homogeneous ideal in $K[x_0, \ldots, x_n]$, i.e.\ it is a set $X \subseteq \pnk$ such that
    \[ X = V_+(I) := \left\{x \in \pnk \ \middle| \ f(x) = 0 \text{ for all } f \in I 
    \text{ homogeneous}\right\},\]
    where $I \subseteq K[x_0 \ldots, x_n]$ is a homogeneous ideal.
\end{defn}

As in the affine case, a projective variety is equivalently described as the vanishing locus of a finite set of homogeneous polynomials. One may define an operator $I_+$ associating to a projective variety $X$ the homogeneous ideal generated by homogeneous polynomials that vanish on $X$.

Projective varieties are precisely the closed sets of a topology on $\pnk$, and a projective variety is irreducible if and only if $I_+(X)$ is prime. Furthermore, projective varieties can be uniquely decomposed into irreducible components.

A particularity of projective space is that it locally looks like affine space. Formally, it is covered by the standard open sets
\[ U_i := D_+(x_i) = \pnk \setminus V_+(x_i).\]
In particular, this means that a subset $Z \subseteq \pnk$ is closed (i.e., a projective variety) if and only if $Z \cap U_i$ is closed in $U_i$ (in the affine Zariski topology) for all $i=0, \ldots, n$. This technique will be called the \emph{affine criterion for closedness}.

\subsubsection{Additional notions}

To complete the foundational toolkit, we introduce several constructions that allow us to manipulate varieties and relate them to linear algebra.

\paragraph{Projectification of vector spaces}
Just as projective space $\pnk$ is derived from the affine space $K^{n+1}$, any finite-dimensional vector space $W$ can be \emph{projectified}. We denote this by $\PP(W)$. Formally, $\PP(W)$ is the set of lines passing through the origin in $W$, or equivalently, the quotient $(W \setminus \{0\}) / K^*$. If $\dim W = n+1$, then $\PP(W) \cong \PP^n$.


\paragraph{Product of varieties}
If $X \subseteq \ank$ and $Y \subseteq \affn{m}{K}$ are affine varieties, their Cartesian product $X \times Y$ carries a natural structure of an affine variety in $\affn{n+m}{K}$. Topologically, the Zariski topology on $X \times Y$ is not the product of the Zariski topologies on $X$ and $Y$; rather, the closed sets are defined by polynomials in the combined variables $x_1, \ldots, x_n, y_1, \ldots, y_m$. The product of two projective varieties can be defined similarly. However it is worth noting that, to give it the structure of a projective variety, it needs to be embedded into a bigger projective space via the so-called \emph{Segre embedding}. Formally, if $X \subseteq \pnk$ and $Y \subseteq \proj{m}{K}$ are two projective varieties, their product $X \times Y \subseteq \pnk \times \proj{m}{K}$ is embedded into $\proj{(n+1)(m+1) - 1}{K}$ via the map $\phi: \pnk \times \proj{m}{K} \rightarrow \proj{(n+1)(m+1) - 1}{K}$ given by 
\[\left(\left[x_0: \cdots: x_n], [y_0: \cdots:y_m\right]\right) \mapsto \left[x_0 y_0: x_0 y_1: \cdots: x_0 y_m: \cdots: x_n y_m\right]\]

\paragraph{Coordinate ring}
Let $X \subseteq \ank$ be an affine variety. The \emph{coordinate ring} of $X$, or \emph{ring of regular functions}, written $k[X]$, is defined as
\[k[X] = \frac{K[x_1, \ldots, x_n]}{I(X)}.\]
The projective version is defined analogously.\\
The quotient allows to identify polynomial functions over $\ank$ that have the same value on $X$, which explains the name of regular functions.

\paragraph{Morphisms of varieties}
For the sake of this project, we will not give the formal definition of a morphism of varieties, but only the intuition behind it. Intuitively, a \emph{morphism} $\phi: X \to Y$ between affine varieties is simply a function whose components are given by ratios of polynomials. That is, if $X \subseteq \affn{n}{K}$ and $Y \subseteq \affn{m}{K}$, then $\phi(x) = (f_1(x)/g_1(x), \ldots, f_m(x)/g_m(x))$ where each $f_i, g_i \in k[X]$ has image in $Y$, and each $g_i$ is nowhere zero on $X$. For projective varieties, morphisms are locally given by ratios of homogeneous polynomials of the same degree. Morphisms are continuous with respect to the Zariski topology.

\subsubsection{Tangent space of a point}\label{subsubsection:tangent}
Just like in differential geometry, there is a way to associate to any point of a variety a tangent space. The techniques to get there, instead of being analytical, will be purely algebraic.
\begin{defn}
    Let $X \subseteq \ank$ be an affine variety. The \emph{coordinate ring} of $X$, or \emph{ring of regular functions} written $k[X]$, is defined as
    \[k[X] = \frac{K[x_1, \ldots, x_n]}{I(X)}.\]
    The projective version is defined analogously.
\end{defn}
The quotient allows to identify polynomial functions over $\ank$ that have the same value on $X$, which explains the name of regular functions.

\begin{defn}
    Let $X \subseteq \ank$ be an affine variety, and $x \in X$ be a point. The \emph{local ring of $X$ at $x$}, denoted $\mathcal{O}_{X,x}$, is the ring of all functions $f = g/h$, where $g, h \in k[X]$ and $h(x) \neq 0$. Since the definition is purely local, the definition in the projective case is done by working on an affine chart.
\end{defn}
It is indeed a local ring, with its maximal ideal $\mathfrak{m}_x$ consisting of functions $f = g/h$ where $g(x) = 0$.

\begin{defn}
    Let $X \subseteq \ank$ be an affine variety, and $x \in X$ be a point. The \emph{tangent space} of $X$ at $x$, denoted $\mathcal{T}{(X)}_x$, can be seen as the space of point derivations, i.e.\ of $K$-linear maps $\delta : \mathcal{O}_{X,x} \to K$ such that 
    \[ \delta\left(fg\right) = \delta(f) g(x) + f(x) \delta(g). \]
    Since the definition is local, the projective version is defined by using an affine chart.
\end{defn}
A point derivation is uniquely determined by its effect on $k[X]$ embedded into $\mathcal{O}_{X,x}$. Indeed, if $f \in k[X]$ is such that $f(x) \neq 0$, then
\[ 0 = \delta\left(1\right) = \delta\left(\frac{f}{f}\right) = \delta(f) \frac{1}{f(x)} + \delta\left(\frac{1}{f}\right) f(x), \]
so that 
\[\delta\left(\frac{1}{f}\right) = - \frac{\delta(f)}{f{(x)}^2}, \]
which is well-defined since $f(x) \neq 0$. It is worth noting that the tangent space is more commonly defined as the dual vector space ${(\mathfrak{m}_x / \mathfrak{m}_x^2)}^*$, but one can show that both vector spaces are isomorphic.

The view of the tangent space as point derivations allows for the definition of \emph{differential of a morphism}. 

\begin{defn}
    Let $X \subseteq \ank$, $Y \subseteq \affn{m}{K}$ be two affine varieties, $x \in X$ a point, and $\phi : X \rightarrow Y$ be a morphism of varieties. The \emph{differential} of $\phi$ at $x$ is the map $d \phi_x : \mathcal{T}{(X)}_x \to \mathcal{T}{(Y)}_{\phi(x)}$ defined via
    \[ d\phi_x(\delta)(f) = \delta\left(f \circ \phi\right) \quad \left(\delta \in \mathcal{T}{(X)}_x, f \in \mathcal{O}_{Y,\phi(x)}\right).\]
\end{defn}

The fact that $f \circ \phi$ belongs to $\mathcal{O}_{X,x}$ is hidden in the definition of a morphism of varieties, so that the definition makes sense.

\subsection{Grassmannians, the Plücker embedding and flag varieties}
For the study of algebraic groups, specific spaces arise naturally, namely flag varieties, which are fundamentally related to Grassmannians. In this section, we will see how Grassmannians and flag spaces can be embedded into projective space by means of the Plücker embedding.

\subsubsection{Reminders in linear algebra}

Let $V$ be an $n$-dimensional vector space over $K$. The \emph{exterior algebra} $\wedge V$ is the quotient of the tensor algebra
\[ \text{T}V := \bigoplus_{k \geq 0} \text{T}^k V = K \oplus V \oplus (V\otimes V) \oplus \ldots \]
by the ideal generated by all $v \otimes v$ for $v \in V$. Elements of $\wedge V$ are therefore finite sums of exterior products $v_1 \wedge v_2 \wedge \cdots \wedge v_k$, where the exterior product is the class of $v_1 \otimes \cdots \otimes v_k$ in $\wedge V$. Using the quotient relation, the exterior product is commutative up to a sign. 

Choosing an ordered basis of V naturally yields a basis of $\wedge V$. Namely, if $v_1, \ldots, v_n$ is a basis of $V$, then the wedge products $v_{i_1} \wedge \ldots v_{i_d}$ (for $i_1 < i_2 < \cdots < i_d$) give a basis of $\wedge^d V$. Consequently, each $\wedge^d V$ is a $\binom{n}{d}$-dimensional $K$-vector space, proving finite dimensionality of $\wedge V$. In particular, any two basis of $V$ have the same exterior product up to a non-zero scalar, since $\wedge^n V$ is $1$-dimensional. 

Note that $\wedge V$ is a finite dimensional graded algebra over $K$, with $\wedge^i V = V \wedge \cdot \wedge V$ the $i$-th wedge product. Because of the anticommutativity of the exterior product, any exterior product of more than $n$ vectors vanishes (due to them being linearly dependent). As a result, we see that $\wedge^d V = 0$ whenever $d > n$.

For any subvector space $W \subseteq V$, note that $\wedge^d W$ is canonically identified with a subspace of $\wedge^d V$ by considering the wedge products $w_1 \wedge \cdots w_d$ in $\wedge^d V$ (using that $w_i \in V$ for all i) instead of in $\wedge^d W$.

The \emph{symmetric algebra} $\Sym{}{V}$ is the quotient of $\textnormal{T}(V)$ by the ideal generated by the elements $v \otimes w - w \otimes v$ ($v,w \in V)$. Its elements are finite sums of arbitrary length tensor products $v_1 \otimes \ldots \otimes v_n$, where the tensor is made commutative by the quotient relation. The tensor product $\otimes$ is replaced by $\odot$ to indicate the commutativity of multiplication in $\Sym{}{V}$
It has a natural structure of graded algebra
\[ \Sym{}{V} = \bigoplus_{k = 0}^{\infty} \Sym{k}{V}, \]
where $\Sym{k}{V}$ consists of all $k$-fold symmetric products, and is called the \emph{$k$-th symmetric power} of $V$. While $\Sym{}{V}$ is infinite dimensional, each $\Sym{k}{V}$ is finite-dimensional, of dimension $\binom{k + n - 1}{n}$, with a basis given by all possible combinations of $k$ vectors from a basis of $V$, up to permutation.

\subsubsection{Embedding the Grassmannians into projective space}

Certain projective varieties will arise naturally in the study of algebraic groups, all built on the notion of Grassmannians.

\begin{defn}
    Let $V$ be a $n$-dimensional vector space over $K$ and let $d \leq n$ be a non-negative integer. The \emph{$d$-th Grassmannian of $V$} $\Gr_d(V)$ is the collection of all $d$-dimensional subspaces of $V$.
\end{defn}

Using the previous discussion about exterior algebras, for any $D \in \Gr_d(V)$, there is a natural way to identify the one-dimensional vector space $\wedge^d D$ as a subspace of $\wedge^d V$. This allows us to define the map

\begin{alignat*}{4}
    \psi: \quad &\Gr_d(V) &\ \to\ &\PP(\wedge^d V)\\
           & \quad D &\ \mapsto\ &\left[w_1 \wedge \cdots \wedge w_d \right],
\end{alignat*}

where $w_1, \ldots, w_d$ is a basis of $D$. This map does not depend on the choice of a basis of $D$, so that it is well-defined. Indeed, let $w_1' , \ldots, w_d'$ be another basis of $D$. Each of its vectors can be written as $w_i' = \sum_{j = 1}^d \lambda_{i,j} w_j$ for some $\lambda_{i,j} \in \CC$. Taking the exterior product, all products containing a $w_i \wedge w_i$ are zero, so that we get
\[ w_1' \wedge \cdots \wedge w_d' = \sum_{j = 1}^d \lambda_{1,j} w_j \wedge \cdots \wedge \sum_{j = 1}^d \lambda_{n,j} w_j = \lambda w_1 \wedge \cdots \wedge w_d, \]
where $\lambda \in \CC$ depends on the $\lambda_{i,j}$, and is non-zero since $w_1', \ldots, w_d'$ is a basis. In particular, taking the projectification yields that $w_1 \wedge \cdots \wedge w_d = w_1' \wedge \cdots \wedge w_d'$ in projective space.

This map is called the  \emph{Plücker embedding}. Its image is in fact closed, allowing us to see Grassmannians as projective varieties.

\begin{theorem}\label{thm:plucker}
    Let $V$ be an $n$-dimensional $K$-vector space, $d \leq n$ a non-negative integer. The map
    \[ \psi: \Gr_d(V) \rightarrow \PP(\wedge^d V) \]
    is injective and its image is closed.
\end{theorem}
To prove this theorem, we need an elementary result about graphs of morphisms.
\begin{lemma}
    Let $f: \affn{n}{K} \to \affn{m}{K}$ be a morphism of varieties, where $f = (f_1, \ldots, f_m)$ with each $f_i\in K(\ank)$. The graph of $f$
    \[\Gamma(f) = \left\{(x_1, \ldots, x_n, f_1({(x_i)}_{i = 1}^n), \ldots, f_m({(x_i)}_{i = 1}^n)) ~\middle|~ (x_1, \ldots, x_n) \in \ank\right\} \subseteq \affn{n+m}{K}\]
    is closed.
\end{lemma}
\begin{proof}
    Note that $\Gamma(f)$ is defined as the vanishing locus of the polynomials $\left\{x_{n + j} - f_j({(x_i)}_{i = 1}^n)\right\}$ in $K[x_1, \ldots, x_{n+m}]$, for $1 \leq j \leq m$. Thus it is closed.
\end{proof}
\begin{proof}[Proof of \cref{thm:plucker}]
    We follow the proof of~\cite[Section~1.8]{humphreys_linear_1975}. Let us first show injectivity. Let $D$, $D'$ be $d$-dimensional subspaces of $V$ such that $\psi(D) = \psi(D')$. Choose a basis $v_1, \ldots, v_n$ of $V$ such that $v_1, \ldots, v_d$ span $D$ and $v_{r + 1}, \ldots, v_{r+d}$ span $D'$ for some $r \geq 0$. Their images via $\psi$ being equal precisely means that $v_1 \wedge \ldots \wedge v_d = \lambda v_{r+1} \wedge \ldots \wedge v_{r+d}$ for some $\lambda \in K^*$. Using the basis of $\wedge^d V$ arising from $v_1, \ldots, v_n$, we see that $v_1 \wedge \ldots \wedge v_d$ is colinear to $v_{r+1} \wedge \ldots \wedge v_{r+d}$ if and only if $r = 0$, i.e. $D=D'$. Thus $\psi$ is injective.

    By the affine criterion for closedness, it suffices to check that $\im\psi$ intersects each member of a standard affine chart of $\PP(\wedge^d V)$ in a closed subset. The extreme cases $d=1$ and $d=n$ are trivial: $\Gr_1(V) = \PP(V)$ and $\Gr_n(V)$ is a single point.

    Fix an ordered basis $v_1, \ldots, v_n$ of $V$.  The associated standard basis of $\wedge^d V$ consists of all elements
    \[
        v_{i_1} \wedge \cdots \wedge v_{i_d} \quad \left(1 \le i_1 < \cdots < i_d \le n\right).
    \]
    For each such multi-index $I = \{i_1,\ldots,i_d\}$, define the affine open set
    \[
        U_I := \left\{[\omega] \in \PP(\wedge^d V) ~\middle|~
            \text{the } v_I\text{-coordinate of }\omega \text{ is nonzero}\right\},
    \]
    i.e.\ the complement of the hyperplane $\{v_{i_1}\wedge\cdots\wedge v_{i_d} = 0\}$.
    These sets $U_I$ cover $\PP(\wedge^d V)$.
    
    Without loss of generality, let us show that the intersection of $\im \psi$ with the chart $U_{1, \ldots, d}$ is closed.
    Consider the chart $U := U_{\{1,\ldots,d\}}$, corresponding to the basis
    element $v_1 \wedge \cdots \wedge v_d$.  Set $D_0 := \vspan{v_1,
    \ldots, v_d}$, the ``reference'' $d$-plane.

    A subspace $D \in \Gr_d(V)$ satisfies $\psi(D) \in U$ if and only if the
    natural projection $\pi : V \twoheadrightarrow D_0$ restricts to an isomorphism
    $\pi|_D : D \xrightarrow{\cong} D_0$.  In particular, $D$ has a unique
    canonical basis of the form
    \begin{equation*}
        \tilde{v}_i \;:=\; v_i + x_i(D), \quad 1 \le i \le d,
        \qquad \text{where } x_i(D) := \sum_{j > d} a_{ij} v_j \in
        \operatorname{span}(v_{d+1},\ldots,v_n),
    \end{equation*}
    given by taking preimages of each $v_i$ by $\pi|_D$.  The
    parameters $a_{ij}$ with $1 \le i \le d$ and $d+1 \le j \le n$ give coordinates
    on this chart; there are $d(n-d)$ of them, matching ${\dim\Gr_d(V) = d(n-d)}$.

    Using the canonical basis, the image of $D$ via the Plücker embedding is
    \begin{align*}
        \psi(D) &= \tilde{v}_1 \wedge \cdots \wedge \tilde{v}_d\\
        &= \bigl(v_1 + x_1(D)\bigr) \wedge \cdots \wedge \bigl(v_d + x_d(D)\bigr)\\
        &= v_1 \wedge \cdots \wedge v_d
        \;+\; \sum_{i = 1}^d v_1 \wedge \cdots \wedge x_i(D) \wedge \cdots \wedge v_d
        \;+\; (\ast),
    \end{align*}
    where $(\ast)$ collects all terms in which two or more of the $v_i$ have
    been replaced by the correction term $x_i(D)$.

    Since $x_i(D) = \sum_{j>d} a_{ij} v_j$, the term
    $v_1 \wedge \cdots \wedge x_i(D) \wedge \cdots \wedge v_d$ expands as
    \[
        \sum_{j>d} a_{ij}\,
        v_1 \wedge \cdots \wedge \widehat{v}_i \wedge \cdots \wedge v_d \wedge v_j,
    \]
    where $\widehat{v}_i$ means $v_i$ is omitted.  In the affine chart $U$, the coefficient of
    the basis element ${v_1 \wedge \cdots \wedge \widehat{v}_i \wedge \cdots \wedge
    v_d \wedge v_j}$ is exactly $\pm a_{ij}$.  Thus the $d(n-d)$ scalars $a_{ij}$
    are recovered as affine coordinates of $\psi(D)$ in $U$.

    Each term in $(\ast)$ involves at least two correction factors $x_i(D)$; its
    coefficient is therefore a polynomial function of the $a_{ij}$.
    Specifically, a term involving corrections at positions $i_1 < i_2$ has a
    coefficient which is a degree-$2$ polynomial in the $a_{ij}$; higher-order terms
    give polynomials of higher degree.  In particular, they do not depend on the choice of a space $D$, only on the choice of some coefficients $a_{ij}$.

    Thus, in the affine chart $U$, the image $\im\psi \cap U$ consists
    precisely of the points with affine coordinates
    \[
        \bigl(\ldots,\, a_{ij},\, \ldots,\, f_k(a_{ij}),\, \ldots\bigr),
    \]
    where the $a_{ij}$ range freely over $K^{d(n-d)}$, and each $f_k$ is an
    explicit polynomial in the $a_{ij}$ arising from the above expansion of $\psi(D)$.

    In other words, $\im\psi \cap U$ is the graph of the regular map
    \[
        \varphi : \affn{d(n-d)}{K} \longrightarrow \affn{N - d(n-d)}{K},
        \quad (a_{ij}) \mapsto (f_k(a_{ij})),
    \]
    where $N = \binom{n}{d} - 1 = \dim\PP(\wedge^d V)$.  The graph of any morphism of
    affine varieties is closed in the product (in the Zariski topology). Therefore
    $\im\psi \cap U$ is closed in $U$.

    Since this holds for every chart $U_I$ in the standard cover, we conclude that $\im\psi$ is closed in $\PP(\wedge^d V)$.
\end{proof}
\noindent Since Grassmannians are projective varieties, it follows that their cartesian product can also be given the structure of a projective variety.

\subsubsection{Flag spaces as projective varieties}

One last notion that will be of particular interest for the study of algebraic groups is that of flags and flag varieties.

\begin{defn}
    Let $V$ be a $K$-vector space. 
    
    A \emph{flag} in $V$ is a chain $0 \subsetneq V_1 \subsetneq \ldots \subsetneq V_k = V$ of proper inclusion of subspaces of $V$. A flag is \emph{full} if $k = \dim{V}$. The \emph{flag variety} of $V$, denoted $\flag{V}$, is the set of all full flags of $V$. One may similarly define \emph{partial flag varieties} of type $(d_1, \ldots, d_k)$, where $0 < d_1 < \cdots < d_k = \dim V$, to be the set of partial flags $0 \subsetneq V_1 \subsetneq \ldots \subsetneq V_k = V$ such that $\dim{V_i} = d_i$.
\end{defn}

It is clear that flag varieties embed into products of Grassmannians. Furthermore, they are defined by closed polynomial equations, making them projective varieties, as is stated in the following proposition.

\begin{theorem}
    The flag variety $\flag{V}$ is closed in $\Gr_1(V) \times \cdots \times \Gr_n(V)$, making it a projective variety.
\end{theorem}
\begin{proof}
    We follow the argument in~\cite[Section~1.8]{humphreys_linear_1975}. Let $n = \dim V$. Consider the product of Grassmannians $\mathcal{X} = \Gr_1(V) \times \Gr_2(V) \times \cdots \times \Gr_{n-1}(V)$. Since each $\Gr_d(V)$ is a projective variety, their product $\mathcal{X}$ is also a projective variety. The flag variety $\flag{V}$ identifies with the subset of $\mathcal{X}$ consisting of chains $V_1 \subset V_2 \subset \cdots \subset V_{n-1}$. To show $\flag{V}$ is closed, it suffices to show that for any $d$, the set of pairs $(D, D') \in \Gr_d(V) \times \Gr_{d+1}(V)$ such that $D \subset D'$ is closed.

    To avoid cumbersome notation, we work locally on affine open subsets. Fix a basis $v_1, \ldots, v_n$ of $V$. Consider the standard affine open subsets $U \subset \PP(\wedge^d V)$ and $U' \subset \PP(\wedge^{d+1} V)$ defined relative to the basis elements $v_1 \wedge \cdots \wedge v_d$ and $v_1 \wedge \cdots \wedge v_{d+1}$ respectively. These sets consist of points where the corresponding coordinate is non-zero. 
    
    If a subspace $D \in \Gr_d(V)$ has its image in $U$, it admits a canonical basis of the form $v_i + x_i(D)$ for $1 \leq i \leq d$, where $x_i(D) = \sum_{j > d} a_{ij} v_j$. Similarly, if $D' \in \Gr_{d+1}(V)$ has its image in $U'$, it admits a canonical basis $v_i + y_i(D')$ for $1 \leq i \leq d+1$, where $y_i(D') = \sum_{j > d+1} b_{ij} v_j$. The coefficients $a_{ij}$ and $b_{ij}$ serve as local coordinates on these affine charts.

    A quick computation with these bases shows that $D \subset D'$ if and only if the basis vectors of $D$ can be expressed as linear combinations of the basis vectors of $D'$. Specifically, this inclusion holds if and only if:
    \[ x_i(D) = y_i(D') + a_{i, d+1}(v_{d+1} + y_{d+1}(D')) \]
    for all $1 \leq i \leq d$. Comparing the coefficients of the basis vectors $v_j$ (for $j > d+1$) on both sides, this condition translates into the polynomial equations:
    \[ a_{ij} = b_{ij} + a_{i, d+1} b_{d+1, j} \]
    for all $1 \leq i \leq d$ and $j > d+1$. Since these are polynomial conditions on the local coordinates $a_{ij}$ and $b_{ij}$, the set of pairs $(D, D')$ satisfying $D \subset D'$ is closed in the affine chart $U \times U'$. As the space is covered by such charts, the incidence relation is closed globally. Thus, $\flag{V}$ is a closed subset of $\mathcal{X}$, making it a projective variety.
\end{proof}
